\section{Model and Trust Boundary}\label{sec:model}
The retirement study of Section~\ref{sec:retire} asks when history can
be discarded; this section and the two that follow define the
verification substrate that makes the discarded history
reconstructable and auditable---the stable identities, generations,
and evidence layouts a typed capsule refers to. These component
results carry their v5.1 scope labels unchanged; they are no longer
the paper's center, and the new ablation
(Table~\ref{tab:retire-evidence}) now bounds when their layout matters.
\subsection{Semantic registry, execution, and display}\label{sec:registry}
Let $\R=(g_0,\ldots,g_{n-1})$ be a nonempty approved registry of unique check
IDs. The candidate identity is
\begin{equation}e=( \text{epoch},\ \text{snapshot},\ \text{checker},\ \text{environment},\ \text{registry}).\end{equation}
The snapshot covers source and protected acceptance inputs. The checker and
environment are explicitly bound; an epoch increases even after an ABA return to
identical source bytes. Only a trusted runner submits final \PASS{}, \FAIL{}, or
\UNKNOWN{} records. Missing, stale, skipped, incompatible, or timed-out outcomes
are not passes.

\paragraph{Immutable identities versus mutable execution state.}
The model splits its state into two disjoint layers.
\emph{Immutable check identities} are the registry $\R$, the candidate
identity $e$, and the acceptance predicate each check declares; they are fixed
at publication and are never rewritten by execution, layout choice, or
learning. \emph{Mutable execution state} is everything a run may change: the
result tokens, their generation counters, epoch-local timestamps, and the
execution permutation $\sigma_e$ that determines when checks run. A record's
identity is its semantic ID bound to $e$; its display position is not part of
that identity. This separation is what makes evidence serialization under
fixed schedules well-posed: the schedule and outcomes are observed inputs, and
only the aggregation layout is chosen. No layout or learner writes execution
state, and no execution mutates the registry or acceptance predicates.

An execution permutation $\sigma_e$ determines when checks run. A layout
$L=(\pi,T)$ independently maps tree positions to semantic registry indices. For
a node covering positional interval $[i,j]$, its semantic coverage is
\begin{equation}G_L(i,j)=\{g_{\pi(i)},\ldots,g_{\pi(j)}\}.\label{eq:coverage}\end{equation}
The stored layout must contain an exact permutation, and every internal node's
children must partition its positional interval without gaps or overlaps. The
task graph is outside this mathematical tree; one worker can discharge many
checks, and one integration check can cover several workers.

The acceptance predicate is independent of both permutations:
\begin{equation}A(e)=\mathbf 1\{\text{every }g_i\in\R\text{ has trusted, matching, current \PASS{} evidence for }e\}.\label{eq:accept}\end{equation}
An observed fail permits rejection before all checks finish. It does not assert
coverage of the unexecuted remainder. Acceptance is relative to the approved
predicates, not proof of all program behavior.

\subsection{Packet and lifecycle costs}\label{sec:packet}
Each internal node of arity $d$ emits a fixed-width packet
\begin{equation}c(d)=H+Sd,\qquad 2\le d\le b=\left\lfloor (C-H)/S\right\rfloor.\label{eq:price}\end{equation}
Here $H=512$, $S=320$, and $C=3072$ bytes, so $b=8$. A mandatory field that
exceeds its slot is rejected, not silently truncated. The cap applies to one
certificate packet, not to all project memory, an entire registry manifest, or
every LLM prompt. A manifest is a referenced control-plane artifact; its measured
size is charged on migration.

Let $W_{e,1},\ldots,W_{e,q_e}$ be the semantic IDs whose final records become
available between successive refreshes. For a strict new candidate epoch, the
observed cost is
\begin{equation}\ell_e(L)=W(L)+\sum_{a=1}^{q_e}\sum_{v\in\mathrm{int}(T)}c(d_v)\,\mathbf 1\{W_{e,a}\cap G_L(v)\ne\varnothing\},
\qquad W(L)=\sum_{v}c(d_v).\label{eq:lcost}\end{equation}
The first term constructs the full initially unknown tree. The CLI batches an
executed prefix into groups of eight, flushing its final partial group. A
successful strict epoch runs the whole registry. This conservative
implementation does not reuse a pass across distinct source snapshots.

We separately study persistent-state refresh traces that omit $W(L)$ at each
tick, as in a cache with correctly specified incremental dependencies. Each
topology installation is still cold and explicitly priced. These traces evaluate
the controller's cost model; they are not the strict candidate-gate timing
experiment. They do not demonstrate safe dependency inference from arbitrary
code.

\subsection{Why more than invalidation or pairwise similarity is needed}\label{sec:why}
A classical result shows that marginal check activation cannot identify the best
aggregation tree under joint invalidation (Appendix~\ref{app:sep}). A further
result shows that even the full initial invalidation distribution cannot identify
optimal progressive verification: early rejection and completion batching change
the update sequence. Both distinctions remain relevant after allowing
permutations.

\begin{proposition}[Pairwise affinity is insufficient for full-wave scoring]\label{prop:pairwise}
There exist two distributions with identical singleton and pairwise co-activity
probabilities, hence identical Jaccard affinities, but different activation
probabilities for the same three-check subtree.
\end{proposition}
\begin{proof}
On IDs $\{0,1,2\}$, let $P$ be uniform on $\varnothing,\{0,1\},\{0,2\},\{1,2\}$
and $Q$ uniform on $\{0\},\{1\},\{2\},\{0,1,2\}$. Every singleton has
probability $1/2$, every pair has joint probability $1/4$, and its Jaccard
affinity is $1/3$. Yet the probability that the triple is hit is $3/4$ under $P$
and $1$ under $Q$. Thus a packet emitted on that interval has different expected
cost. \qed
\end{proof}
This elementary parity construction is not a new fact about higher-order
statistics. Its implication for the compiler is specific: a pairwise clustering
objective cannot substitute for scoring recorded completion hyperedges. It does
not, by itself, prove that these two distributions have different globally
optimal layouts.
